%==============================================================================
% Moumita Dey -- Cover Letter
%
% Companion to cv.tex, from the same "ModernCV and Cover Letter" template.
% Keep the preamble in sync with cv.tex so both documents look identical.
%
% Compile with:  pdflatex cover_letter.tex
%==============================================================================

\documentclass[11pt,a4paper,sans]{moderncv}

%------------------------------------------------------------------------------
% MODERNCV THEME  (must match cv.tex)
%------------------------------------------------------------------------------
\moderncvstyle{classic}
\moderncvcolor{blue}

\definecolor{cvblue}{RGB}{21,76,143}
\colorlet{color1}{cvblue}
\colorlet{color2}{cvblue}

%------------------------------------------------------------------------------
% PACKAGES
%------------------------------------------------------------------------------
\usepackage[utf8]{inputenc}
\usepackage[T1]{fontenc}
\usepackage{lmodern}
\usepackage[english]{babel}
% Matches cv.tex, where 0.82 is needed to keep the Systems section whole.
\usepackage[scale=0.82]{geometry}

\setlength{\hintscolumnwidth}{2.9cm}

%------------------------------------------------------------------------------
% PERSONAL DATA  (must match cv.tex)
%------------------------------------------------------------------------------
\name{Moumita}{Dey}
% \mbox keeps "AML & COMPLIANCE" together, so the tagline wraps after the
% second bar rather than stranding "COMPLIANCE" on its own line.
\title{BANKING \textbar{} KYC \textbar{} \mbox{AML \& COMPLIANCE}}
\address{Zurich, Switzerland}{}{}
\phone[mobile]{+41~76~711~16~39}
\email{reachmoumita.93@gmail.com}
\extrainfo{Swiss Permit B}
\photo[70pt][0.4pt]{resources/photo}

%==============================================================================
\begin{document}

%------------------------------------------------------------------------------
% RECIPIENT  -- replace with the details of the role you are applying for
%------------------------------------------------------------------------------
\recipient{Hiring Manager}{Company Name\\Street and Number\\8000 Zurich}
\date{\today}
\opening{Dear Hiring Manager,}
\closing{Yours sincerely,}
\enclosure[Enclosed]{curriculum vit\ae{}}

\makelettertitle

I am writing to apply for the \emph{[position title]} role at \emph{[company
name]}. I am a banking and wealth management professional with nearly nine years
of experience across client onboarding, KYC and customer due diligence, AML
escalations, fraud and card operations and the servicing of High Net Worth and
affluent clients. I hold a Swiss B permit and am based in Zurich, so I require
no employer sponsorship.

In my most recent role as Chief Manager at IndusInd Bank, I managed
relationships with affluent and High Net Worth clients while supporting KYC
verification, remediation and periodic refresh activities. That work involved
reviewing customer information and supporting documentation, resolving gaps in
Source of Funds and Source of Wealth files, and supporting PEP and sanctions
screening, UBO verification and FATCA/CRS requirements. Where customer activity
or documentation warranted further review, I escalated it to the appropriate
compliance and AML teams. Across earlier roles at HDFC Bank and at Axis
Bank, where I rose from Assistant Manager to Deputy Manager, I built the same
strengths over a broader product range covering deposits, lending, investments,
insurance, foreign exchange and digital banking.

I have recently strengthened my regulatory grounding for the Swiss market
through the SAQ-recognised AMLA/AMLO-FINMA 2025 training on the prevention of
money laundering, and through the Basel Institute on Governance course on
combating terrorism financing, alongside further certifications in risk
management and AML/KYC for banking and insurance.

What I would bring to \emph{[company name]} is a client-facing compliance
perspective: I am used to obtaining difficult documentation directly from
customers and to coordinating between Operations, Compliance and Fraud teams to
close cases properly rather than quickly. I would welcome the chance to discuss
how that experience fits your team.

\makeletterclosing

\end{document}
